% Generated draft: requires mathematical and source review.
\documentclass[11pt]{article}
\usepackage[margin=25mm]{geometry}
\usepackage{fontspec}
\setmainfont{Latin Modern Roman}
\usepackage{amsmath,amssymb,mathtools,unicode-math}
\setmathfont{Latin Modern Math}
\usepackage{enumitem,needspace,xurl,hyperref}
\setlength{\parindent}{0pt}
\setlength{\parskip}{5pt}
\setlength{\emergencystretch}{3em}
\newcommand{\N}{\mathbb{N}}
\newcommand{\Z}{\mathbb{Z}}
\newcommand{\Q}{\mathbb{Q}}
\newcommand{\R}{\mathbb{R}}
\newcommand{\C}{\mathbb{C}}
\newcommand{\F}{\mathbb{F}}
\newcommand{\A}{\mathbb{A}}
\newcommand{\PP}{\mathbb{P}}
\newcommand{\E}{\mathbb{E}}
\newcommand{\eps}{\varepsilon}
\newcommand{\dd}{\,\mathrm{d}}
\begin{document}
\section*{UnsolvedMath ID: 30002437}
\textbf{Source problem code:} catalogue entry
\subsection{Definitions and mathematical statement}
For a positive integer parameter \(n\), define
\[
SP(n):=\#\bigl\{(a,b,c)\in\Z_{>0}^{3}:abc+a+b+c=n\bigr\}.
\]
Here solutions are interpreted as ordered triples. The supplied source does not explicitly specify whether permutations are counted separately, nor does it explicitly state the domain of \(n\); the positive integer domain and ordered-triple convention used here require confirmation against the source.

The proposed bound is
\[
\forall\varepsilon\in\R_{>0}\ \exists C(\varepsilon)\in\R_{>0}\ \forall n\in\Z_{>0},\qquad
SP(n)<C(\varepsilon)n^{\varepsilon}.
\]
The constant may depend on \(\varepsilon\) but must be independent of \(n\). This is a question asking whether the displayed assertion holds, not an asserted theorem.
\subsection{Short English statement}
Does the number of positive integer triples satisfying \(abc+a+b+c=n\) admit, for every \(\varepsilon>0\), a uniform upper bound strictly smaller than \(C(\varepsilon)n^{\varepsilon}\)? The parameter domain and counting convention require source confirmation.
\subsection{Sources}
\begin{itemize}
\item Local UnsolvedMath registry, \texttt{data/unsolvedmath/problems.json}, permanent ID \(30002437\), title: ``Subpolynomial Bounds for an Additive Multiplicative Equation.'' Exact source problem code: \texttt{catalogue entry}. Collection: Oberwolfach Reports Open Problems. No individual proposer is supplied; the proposed year is recorded as 2013.
\item Supplied report citation: \emph{Analytic Number Theory (2014)}, Oberwolfach Reports, pp.\ 2963--3037, DOI 10.4171/owr/2013/51. Supplied source URL and literature reference labeled ``original OWR report'': \url{https://doi.org/10.4171/owr/2013/51}. The supplied background describes these reports as documenting mathematical workshops and their open-problem discussions. The citation's 2014 year differs from the recorded proposed year 2013 and the DOI's 2013 component; these details are retained without reconciliation.
\item Literature assessment dated exactly \texttt{2026-08-21}: the supplied curation labels the problem open and states that the sum-product bound remains conjectural in the authoritative source, with no verified primary resolution located. The sole listed checked source is the original OWR report linked above. This is supplied dated literature triage, not an independently verified assessment; time-sensitive status and cited claims require independent verification.
\item The registry labels the statement \texttt{corrected\_verified} and explains that its mathematical content was unchanged while notation, typography, and wording were normalized into a self-contained publication form. This supplied verification label is not independent mathematical review of this draft.
\end{itemize}
\end{document}
